\documentclass[12pt]{article}
\usepackage[utf8]{inputenc}
\usepackage[T1]{fontenc}
\usepackage[brazilian]{babel}
\usepackage{geometry}
\usepackage{amsmath}
\usepackage{listings}
\usepackage{xcolor}
\usepackage{fancyhdr}
\usepackage{titlesec}
\usepackage{enumitem}

\geometry{margin=1in}
\setlength{\headheight}{15pt}
\setlength{\parindent}{0pt}
\setlength{\parskip}{4pt}

\definecolor{primary}{HTML}{1F4E79}
\definecolor{secondary}{HTML}{0B6E99}
\definecolor{lightbg}{HTML}{F4F8FC}
\definecolor{codebg}{HTML}{F7F7F7}
\definecolor{codeframe}{HTML}{D3D3D3}

\titleformat{\section}
  {\Large\bfseries\color{primary}}
  {}
  {0pt}
  {}

\titleformat{\subsection}
  {\large\bfseries\color{secondary}}
  {}
  {0pt}
  {}

\setlist[itemize]{leftmargin=1.4em,itemsep=2pt,topsep=3pt}
\setlist[enumerate]{leftmargin=1.6em,itemsep=2pt,topsep=3pt}

\pagestyle{fancy}
\fancyhf{}
\lhead{\textbf{Programação em Python}}
% \chead{Aula 03: Operadores}
\rhead{Prof. Msc. Nator Junior Carvalho da Costa}
\cfoot{\thepage}
\renewcommand{\headrulewidth}{0.4pt}

\lstset{
    language=Python,
    basicstyle=\ttfamily\small,
    backgroundcolor=\color{codebg},
    frame=single,
    rulecolor=\color{codeframe},
    keywordstyle=\color{blue},
    commentstyle=\color{gray},
    stringstyle=\color{red},
    breaklines=true,
    showstringspaces=false,
    tabsize=4
}

\begin{document}

\begin{center}
{\LARGE \textbf{Lista de Atividades - Aula 03}}\\[4pt]
{\large Operadores Aritméticos, Relacionais e Lógicos}\\[10pt]
\textbf{Disciplina:} Lógica de Programação em Python\\
\textbf{Professor:} Prof. Msc. Nator Junior Carvalho da Costa\\
\textbf{Semestre:} 2026.1
\end{center}

\vspace{8pt}
\colorbox{lightbg}{%
\parbox{\dimexpr\linewidth-2\fboxsep\relax}{%
\textbf{Instruções}\\
Resolva as atividades abaixo em Python. Considere a ordem dos operadores:
\begin{itemize}
    \item Aritméticos
    \item Relacionais
    \item Lógicos
\end{itemize}
\textbf{Importante:} em Python, potência é representada por \texttt{**}.
}}

\section*{Parte 1: Resolva as Expressões}
Para cada expressão, indique o resultado final (\texttt{True} ou \texttt{False}) e mostre o passo a passo.

\begin{enumerate}
    \item \texttt{2 * 2 < 8 + 3}
    \item \texttt{4 * 2 - 1 > 3 * 1 * 9}
    \item \texttt{(5 - 4 / 2 <= 10 + 1) or (6 * 2 - 3 != 5)}
    \item \texttt{(8 / 2 == 3 * 2) and (9 * 9 - 4 >= 8)}
    \item \texttt{(10 / 2 > 1) or (5 == 5)}
    \item \texttt{(7 + 3 > 8) or (2 ** 3 == 8)}
    \item \texttt{(15 / 3 - 2 <= 2 ** 2) and (6 + 4 == 10)}
    \item \texttt{9 - 3 * 2 + 4 > 5}
    \item \texttt{(5 ** 2 >= 20 + 5) and (7 - 2 != 4)}
    \item \texttt{3 * 3 + 1 < 4 ** 2}
    \item \texttt{(2 + 6 / 2 == 5) or (8 - 4 > 1)}
    \item \texttt{(4 ** 2 / 2 == 8) and (10 - 7 <= 5 - 2)}
\end{enumerate}


\end{document}
